%=========== ΛΥΣΗ - 28ο ΘΕΜΑ Α ===========
\providecommand{\theoriaref}[3]{%
\begin{center}
\fcolorbox{maincolor!70!black}{maincolor!8}{%
\begin{minipage}{.88\linewidth}
\textcolor{maincolor!80!black}{\kerkissans{\textbf{#1~\ref{#2}}}}\\[1mm]
#3
\end{minipage}}
\end{center}}
\providecommand{\orismosref}[1]{%
\theoriaref{Ορισμός}{#1}{Η απάντηση δίνεται από τον αντίστοιχο ορισμό.}}
\providecommand{\apodeixiref}[1]{%
\theoriaref{Θεώρημα}{#1}{Η ζητούμενη απόδειξη δίνεται στην αντίστοιχη απόδειξη.}}

\begin{thema}{Α}
\begin{erwthma}
\item \apodeixiref{thewrhma:paragwgos_dynamhs}

\item Για τη συνέχεια μιας συνάρτησης έχουμε:
\orismosref{orismos:synechhs_synarthsh}
Για τη συνέχεια σε κλειστό διάστημα έχουμε:
\orismosref{orismos:synechhs_se_diasthma}

\item \orismosref{orismos:arxikh_synarthsh}

\item Οι χαρακτηρισμοί των προτάσεων είναι:
\begin{alist}
\item \textbf{Σωστό}. Αν $F,G$ είναι δύο παράγουσες της ίδιας $f$ στο $\Delta$, τότε από το θεώρημα αρχικής συνάρτησης διαφέρουν κατά σταθερά.
\item \textbf{Σωστό}. Για κάθε $x\in\mathbb{R}$ ισχύει
\[ \sqrt[3]{x^4}=\left(\sqrt[3]{x}\right)^4=x^{\frac43}, \]
άρα οι δύο συναρτήσεις είναι ίσες.
\item \textbf{Σωστό}. Οι τριγωνομετρικές συναρτήσεις είναι συνεχείς σε κάθε σημείο του πεδίου ορισμού τους.
\item \textbf{Λάθος}. Η ύπαρξη του $\lim\limits_{x\to x_0}f(x)$ δεν εξαρτάται από την τιμή $f(x_0)$, αλλά από τις τιμές της $f$ κοντά στο $x_0$.
\item \textbf{Λάθος}. Για $x>0$ είναι
\[ (x^x)'=\left(e^{x\ln{x}}\right)'=x^x(\ln{x}+1), \]
όχι $x\cdot x^{x-1}$.
\end{alist}

\item Σωστή απάντηση είναι η \textbf{γ}. Το βασικό τριγωνομετρικό όριο είναι
\[ \lim_{x\to0}\dfrac{\hm{x}}{x}=1. \]
\end{erwthma}
\end{thema}
%=========== ΤΕΛΟΣ ΛΥΣΗΣ - 28ο ΘΕΜΑ Α ===========
